%% ma: L11-B04-0012
%% lop: 11
%% chuong: 1
%% bai: 4
%% dang: 11.4.5
%% loai: TN4
%% muc: NB
%% dap_an: "A"
%% nguon: "de-cuong-giua-ky-1-toan-11-nam-2026-2027-truong-thpt-hoang-van-thu-ha-noi-2-8.pdf (M08, Đề cương giữa kì 1 Toán 11 THPT Hoàng Văn Thụ, Hà Nội), Đề minh họa, Phần I Câu 3"
%% kien_thuc: "Bài 4 (phương trình lượng giác cơ bản)"
%% trang_thai: chinh-thuc
%% ghi_chu: "Đã duyệt 29/10/2026."
\begin{ex}
Nghiệm của phương trình $\tan x+\sqrt3=0$ là
\choice{\True $x=-\dfrac{\pi}{3}+k\pi$}{$x=-\dfrac{\pi}{6}+k\pi$}{$x=\dfrac{\pi}{3}+k2\pi$}{$x=\dfrac{\pi}{6}+k\pi$}
\loigiai{$\tan x=-\sqrt3=\tan\left(-\dfrac\pi3\right)\Leftrightarrow x=-\dfrac\pi3+k\pi$ $(k\in\mathbb Z)$.}
\end{ex}
